\section{Genuine finite rehomogenization over the original analytic tangent ring}

Let $B=\mathcal O_{\mathbb C^d,0}$ and $P=B[z_1,\ldots,z_t]$.
Fix an actual normal pivot $j$ and actual complex chart directions $b_i$.
The polynomial algebra map $\delta$ fixes every entire coefficient in $B$,
sends $z_j$ to $1$ and sends $z_i$ to $b_i+z_i$ for $i\ne j$.
Constants $b_i$ here are their genuine constant analytic germs.

\begin{theorem}[Literal finite relative rehomogenization]
\label{hir:actual-relative-finite-polynomial-rehomogenization}
\uses{hir:actual-polynomial-rehomogenization}
\lean{Hironaka.HironakaActualRelativePolynomialRehomogenization.relativePolynomialRehomogenization,Hironaka.HironakaActualRelativePolynomialRehomogenization.relativePolynomialRehomogenization_homogeneous,Hironaka.HironakaActualRelativePolynomialRehomogenization.relativePolynomialDehomogenization_rehomogenization}
\leanok
For a literal $q\in P$ and $m\ge\deg q$, the finite polynomial
\[
 R_m(q)=\sum_{\alpha\in\operatorname{supp}q}
 z_j^{m-|\alpha|}q_\alpha\prod_i(z_i-b_i z_j)^{\alpha_i}
\]
is genuinely homogeneous of normal degree $m$. If $q$ has no pivot
variable, then $\delta R_m(q)=q$. Every coefficient is the unchanged
original analytic tangent germ; no formal homogeneous lift is input.
\end{theorem}
\begin{proof}
For each supported exponent, $|\alpha|\le\deg q\le m$. Each
$z_i-b_i z_j$ is homogeneous of normal degree one over the entire
coefficient ring. Its indicated product has degree $|\alpha|$ and
multiplying by $z_j^{m-|\alpha|}$ gives degree $m$. The actual finite
sum remains homogeneous of that degree. Under $\delta$, every pivot
power becomes $1$. For $i\ne j$, the linear expression becomes
$(b_i+z_i)-b_i=z_i$. A supported exponent of the assumed pivot-free
$q$ has $\alpha_j=0$, so its pivot factor is absent. Thus each entire
summand becomes the original monomial $q_\alpha z^\alpha$.
The finite support expansion of the actual polynomial $q$ proves the
whole dehomogenization identity.
\end{proof}

\begin{theorem}[Actual relative homogeneous dehomogenization is injective]
\label{hir:actual-relative-homogeneous-dehomogenization-injective}
\uses{hir:actual-relative-polynomial-formal-taylor-coefficient,hir:actual-normal-polynomial-blowup-factorization,hir:actual-formal-blowup-all-directions-whole-power-reflection,hir:actual-analytic-blowup-exceptional-quotient-formal-dictionary}
\lean{Hironaka.HironakaActualRelativePolynomialRehomogenization.relativePolynomialDehomogenization_eval,Hironaka.HironakaActualRelativePolynomialRehomogenization.relativePolynomialDehomogenization_homogeneous_injective}
\leanok
Literal original normal-polynomial evaluation of $\delta P$ equals
the actual dehomogenized normal-polynomial analytic germ, retaining all
entire tangent coefficients. If two actual polynomials in $P$ are
homogeneous of the same normal degree and have equal $\delta$ images,
then the whole original polynomials are equal.
\end{theorem}
\begin{proof}
Constants in $B$ are retained by both actual polynomial evaluations.
For the pivot and every other normal variable the two evaluations have
exactly the displayed actual coordinate-germ values. Polynomial
induction by constants, addition and multiplication by a variable
therefore proves the whole evaluation identity.

Original normal-polynomial evaluation is injective: its actual whole
Taylor normal coefficient at any exponent is the actual Taylor series
of that polynomial coefficient, and the actual tangent Taylor map is
injective. Now suppose $P,Q$ have the same normal degree $D$ and
$\delta P=\delta Q$. Genuine normal-homogeneous blowup factorization
gives $\pi^*\operatorname{ev}(P)=e^D\operatorname{ev}(\delta P)$
and the same identity for $Q$. These actual analytic pullbacks are
therefore equal. Their actual whole Taylor series are the genuine
formal blowup images of the respective actual Taylor series. The
proved actual formal blowup is injective in every original direction,
so those original Taylor series agree. Actual Taylor injectivity and
then normal-polynomial evaluation injectivity give $P=Q$ as entire
polynomials over the original analytic coefficient ring.
\end{proof}

\begin{theorem}[Every finite pivot-free relative chart relation has an internally constructed whole homogeneous source relation]
\label{hir:actual-relative-finite-polynomial-relation-rehomogenization}
\uses{hir:actual-relative-finite-polynomial-rehomogenization,hir:actual-relative-homogeneous-dehomogenization-injective}
\lean{Hironaka.HironakaActualRelativePolynomialRehomogenization.exists_relative_homogeneous_polynomial_relation_lift}
\leanok
Let a finite family $f_i\in P$ be homogeneous of degrees $w_i$, and
let pivot-free $q_i\in P$ satisfy the whole polynomial relation
$\sum_iq_i\delta f_i=0$. There are one actual integer $D\ge w_i$
and actual $P_i\in P$ such that $P_i$ is homogeneous of degree
$D-w_i$, $\delta P_i=q_i$, and $\sum_iP_if_i=0$ as a whole
polynomial identity. No chosen degree bound or desired source relation
is an input.
\end{theorem}
\begin{proof}
Choose the genuine finite integer $D=\sum_i(w_i+\deg q_i)$.
Nonnegativity gives $w_i+\deg q_i\le D$ and hence
$\deg q_i\le D-w_i$ for every index. Apply the first theorem to
$P_i=R_{D-w_i}(q_i)$. It gives the asserted homogeneity and exact
$\delta P_i=q_i$. Each product $P_if_i$ is homogeneous of degree
$(D-w_i)+w_i=D$, and so is their finite sum. Its dehomogenization is
$\sum_iq_i\delta f_i=0$ by genuine polynomial homomorphism laws.
The preceding homogeneous injectivity, comparing with the zero
polynomial of degree $D$, therefore makes the whole source sum zero.
All original analytic tangent coefficients have been retained throughout.
\end{proof}

This proves a whole actual relative polynomial relation lift. Analytic
normalization and exceptional regularity require their separate genuine
same-family analytic syzygy argument; no current Hilbert--Samuel or
uniform-history criterion follows merely from this polynomial result.
